%----------------------------------------------------------------------------------------
%    PACKAGES AND THEMES
%----------------------------------------------------------------------------------------

\documentclass[aspectratio=169,xcolor=dvipsnames]{beamer}
\usetheme{SimplePlus}

\usepackage{hyperref}
\usepackage{graphicx} % Allows including images
\usepackage{booktabs} % Allows the use of \toprule, \midrule and \bottomrule in tables
\usepackage{svg} % Added for SVG support
\usepackage{datetime}
\usepackage{bookmark} % Add bookmark package to avoid rerun warnings

% Command to set the date for individual frames
\newcommand{\framedate}[1]{\def\currentframedate{#1}}
% Default frame date
\framedate{\today}

% Define custom footer with logo and module name
\makeatletter
\setbeamertemplate{footline}{%
  \ifnum\c@framenumber>1\relax% Only show footer on non-title pages
    \begin{beamercolorbox}[ht=3.5ex,dp=2ex,leftskip=0.5em,rightskip=0.5em]{footline}%
      \includesvg[height=4ex]{fig/hm_logo.svg}\quad DLVC SoSe25\hfill\currentframedate\quad \insertframenumber{} / \inserttotalframenumber \;%
    \end{beamercolorbox}%
  \fi%
}
\makeatother

% Custom title page with logo
\setbeamertemplate{title page}{%
  \begin{center}
    \includesvg[height=1.5cm]{fig/hm_logo.svg}\\[1em]
    {\usebeamerfont{title}\usebeamercolor[fg]{title}\inserttitle\par}%
    \ifx\insertsubtitle\@empty%
    \else%
      \vskip0.25em%
      {\usebeamerfont{subtitle}\usebeamercolor[fg]{subtitle}\insertsubtitle\par}%
    \fi%
    \vskip1em%
    {\usebeamerfont{author}\insertauthor\par}%
    \vskip0.5em%
    {\usebeamerfont{institute}\insertinstitute\par}%
    \vskip0.5em%
    {\usebeamerfont{date}\insertdate\par}%
  \end{center}
}

%----------------------------------------------------------------------------------------
%    TITLE PAGE
%----------------------------------------------------------------------------------------

\title{Update Meetings}
\subtitle{3D Spatial Scene Description for Blind People}
\author{Lukas Ederer \\ Tobias Lewald \\ Jan Duchscherer}
\institute{Fakult\"at f\"ur Informatik und Mathematik\\Hochschule M\"unchen}
\date{Summer Semester 2025}

%----------------------------------------------------------------------------------------
%    PRESENTATION SLIDES
%----------------------------------------------------------------------------------------

\begin{document}

\begin{frame}
    % Print the title page as the first slide
    \titlepage
\end{frame}

% Overview slide
\framedate{Overview - Summer 2025}
\begin{frame}{Update Meetings --- Overview}
\begin{table}[]
\centering
\begin{tabular}{ll}
\toprule
\textbf{\#} & \textbf{Date} \\
\midrule
1 & 27.03.2025 \\
2 & 01.04.2025 \\
3 & \ldots \\
\bottomrule
\end{tabular}
\end{table}
\end{frame}

% Update Meeting 1 – Tasks
\framedate{27.03.2025}
\begin{frame}{Update Meeting 1 --- Tasks Worked On}
\begin{itemize}
  \item Play with Gemini Spatial understanding and get initial idea of its capabilities and limitations.
  \item Get general understanding of the project goals.
  \item Implement basic voice to text via OpenAI Whisper (needed for some other task).
  \item Reflect upon questions for clarification.
  \item Reflect general ideas for the holistic project.\\
  \( \rightarrow \) Dynamic Semantic Scene Graph, SLAM, 3D native object detection, etc.
\end{itemize}
\end{frame}

% Update Meeting 1 – Problems
\framedate{27.03.2025}
\begin{frame}{Update Meeting 1 --- Problems Encountered}
\begin{itemize}
  \item How does the Data look like? What features are available?
  \item What models are currently being used?
  \item Wie sieht die Entwicklungsumgebung aus (Ger\"ate, iPhone notwendig?)
  \item Entwicklung auf nicht Apple-Ger\"aten.
  \item What parts of the application should finally run on the iPhone?
  \item API Costs: Who will pay? \(\rightarrow \) \textit{SmartAIs}
  \item Initial ambiguities w.r.t.\ project goals.
  \item Gemini 2.0 flashes 3D OOBBs do not seem reliable.
\end{itemize}
\end{frame}

% Update Meeting 1 – Plans
\framedate{27.03.2025}
\begin{frame}{Update Meeting 1 --- Plans for Next Week}
\begin{itemize}
    \item Research on one-shot end to end multimodal detection models.
    \item Await Data from Yamen and do initial exploration.
    \item Explore data structures (StrayScanner, Record3D)
    \item Get Gemini 2.0 flash via API running.
\end{itemize}

\end{frame}

% Placeholder for next meetings
\framedate{01.04.2025}
\begin{frame}{Update Meeting 2}
[Date]
\end{frame}

\framedate{01.04.2025}
\begin{frame}{Update Meeting 2 --- Tasks Worked On}
\begin{itemize}
  \item \textbf{SpatialBot}: VLM fine-tuned on \textbf{RGB-D} data for better 3D understanding~\cite{cai2025spatialbotprecisespatialunderstanding}.
  \begin{itemize}
  \item \( \oplus \) approx. 3B param model.
  \item \( \oplus \) Can explicitly query depth values.
  \item \( \ominus \) no support for multiple video frames.
  \item \( \ominus \) doesn't use camera poses, operates on single images.
  %   \item \textit{SpatialBench} for evaluating VLMs spatial understanding.
  %   \item \textit{SpatialQA}: Question answering dataset given RGB-D data.
  \end{itemize}
  \item \textbf{Video-3D LLM}: 3D spatial understanding - depth information \( \rightarrow \) 3D Pos-Encodings.
  \begin{itemize}
    \item \( \oplus \) multiple video frames + camera poses as input.
    \item \( \ominus \) no explicit detections or querying of depth values.
  \end{itemize}
\end{itemize}
\end{frame}

\framedate{01.04.2025}
\begin{frame}{Update Meeting 2 --- Problems Encountered}
\begin{itemize}
  \item Ambiguities w.r.t.\ project stages.
  \begin{itemize}
    \item Improved visual understanding incorporating depth data?
    \item Route description via dynamic scene graph?
  \end{itemize}
  \item No literature on route description based on RGB-D data.
  \item Focus on downstream application using Gemini Spatial Understanding?
\end{itemize}
\end{frame}

\framedate{01.04.2025}
\begin{frame}{Update Meeting 2 --- Plans for Next Week}
  \begin{itemize}
    \item Combination of depth data w/ detections.
    \item Grobe Überlegungen zur Modellarchitektur.
    \item Explore potential network architectures.
\end{itemize}
\end{frame}

\end{document}